% 阶段研究汇总模板：复制为 NN-主题.tex 后填写，主题按内容定。
% 要求：正式论文风格，逐节引用文献（\cite 引用键须在 references.bib 中存在）。
\documentclass[UTF8,11pt,a4paper]{ctexart}
\usepackage[margin=2.6cm]{geometry}
\usepackage{amsmath,amssymb}
\usepackage{booktabs}
\usepackage[colorlinks=true,linkcolor=blue,citecolor=blue,urlcolor=blue]{hyperref}
\usepackage[backend=biber,style=numeric,sorting=none,doi=true,url=false,eprint=true]{biblatex}
\renewcommand*{\bibfont}{\small}
\setlength{\bibitemsep}{0pt}
\addbibresource{../../literature/references.bib}

\title{第 N 阶段：阶段主题}
\author{Quantum Optics Thesis 工作笔记}
\date{YYYY 年 MM 月 DD 日}

\begin{document}
\maketitle

% 摘要：一段话概括问题、方法、主要结果与证据等级。
\begin{abstract}
本阶段围绕……展开，得到……（注明是数值诊断、数值强支持还是局部候选）。
\end{abstract}

\section{引言}
% 正式论文式文献综述：背景、近年前人工作（逐篇 \cite）、本阶段问题的增量与动机。

\section{模型与方法}
% 哈密顿量、自由能泛函、判据、数值方法（\cite 方法先例，指向 theory/ 与 results/runs/）。

\section{结果}
% 分小节叙述；关键数值附误差边界与证据等级，不把诊断写成严格界。

\section{讨论}
% 与先例的对比、误差与适用范围、可证伪条件。

\section{结论}
% 能说什么、不能说什么、下一步。

\printbibliography[title={参考文献}]
\end{document}
